% !PRINTED 335
\weberpage{335}
%<<CONT
keit von vornherein zu schätzen, als die Abnahme der Differenzen
$\varDelta, \varDelta', \varDelta''\ldots$ Ist $\varDelta'$ so klein,
dass es ausserhalb der Grenzen der beabsichtigten Genauigkeit fällt, so
giebt die Berücksichtigung der ersten Differenz ein genaues Resultat.
Die hier aus einander gesetzte Vorschrift zur Wurzelberechnung wird die
\gesperrt{Regula falsi} genannt.

Die einfachste, für die erste Annäherung geeignete, Form dieser
Vorschrift ist die:

Liegt zwischen $\alpha$ und $\beta$ eine Wurzel $x$ der Gleichung
$f(x)=0$ und ist $f(\alpha) = -a, f(\beta) = b$, so ist
\weq{9}{ x = \frac{b\alpha + a\beta}{a+b} }
ein genäherter Werth von $x$. Dies ist nur eine andere Schreibweise für
die Formel (7).

\begin{center}
§. 109.\\
\gesperrt{Die Newton'sche Näherungsmethode.}
\end{center}

Eine Methode, die zur genäherten Berechnung der Wurzeln einer Gleichung
meist besser ist als die Interpolation, rührt von \gesperrt{Newton} her
und wurde von \gesperrt{Fourier} ausgebildet und genauer untersucht.

Die Methode besteht einfach in Folgendem. Es sei
\weq{1}{ f(x) = 0 }
die aufzulösende Gleichung und es sei ein Werth $x = \alpha$ gefunden,
den man als eine gewisse Annäherung an eine Wurzel betrachten kann. Wir
setzen in (1)
\[ x = \alpha + h, \]
und erhalten
\weq{2}{ f(\alpha+h) = f(\alpha) + hf'(\alpha) +
  \frac{h^2}{1.2}f''(\alpha) + \cdots }
Wenn man nun $h$ aus der Gleichung bestimmt
\weq{3}{ f(\alpha) + hf'(\alpha) = 0, }
was voraussetzt, dass $f'(\alpha)$ von Null verschieden ist, so wird
\[ f(\alpha+h) = \frac{h^2}{1.2}f''(\alpha) + \cdots, \]
und wird also, wenn $h$ eine kleine Zahl ist, da nur das Quadrat
%CONT>>
